\section{L3 OpenMP}

\subsection{LZ1 Introduction \& Model}

\mult{2}

\begin{definition}{What OpenMP Is}
    Three elements:
    \begin{itemize}
        \item runtime library (implementation-dependent)
        \item compiler directives \texttt{\#pragma} (ignored if unsupported)
        \item environment variables (change behaviour at run)
    \end{itemize}
\end{definition}

\begin{concept}{Execution Model}
    \begin{itemize}
        \item shared-memory; Fork-Join + SPMD
        \item a directive opens a \emph{team} of threads
        \item \important{implicit barrier} at end of every parallel region
    \end{itemize}
\end{concept}

\multend

\begin{KR}{Recipe: Build a Parallel Region}
    \paragraph{Open a team} \texttt{\#pragma omp parallel} \{ \ldots \}.
    \paragraph{Share work} add \texttt{omp for} (else every thread repeats the block).
    \paragraph{Mind the barrier} threads join at the implicit barrier (drop it with \texttt{nowait}).
\end{KR}

\begin{example2}{Ex.\ LZ1 Components \& barrier}
    Name the three components of OpenMP. What happens at the end of a parallel region?
     \\ \solution{Solution: }
    \important{Runtime library}, \important{compiler directives}, \important{environment variables}. At the region's end there is an \important{implicit barrier} all threads join before the master continues.
\end{example2}

\subsection{LZ2 Data Scope \& Worksharing}

\mult{2}

\begin{concept}{parallel vs parallel for}
    \begin{itemize}
        \item \texttt{parallel} alone: each thread runs the \emph{whole} block (replicate)
        \item \texttt{parallel for}: iterations \emph{split} (worksharing)
    \end{itemize}
\end{concept}

\begin{definition}{Scope Clauses}
    \begin{itemize}
        \item \texttt{shared}: one copy (race risk)
        \item \texttt{private}: per-thread, uninit, discarded
        \item \texttt{firstprivate}/\texttt{lastprivate}: seed / copy out
        \item \texttt{reduction(op:v)}: private partials, combined
    \end{itemize}
\end{definition}

\multend

\begin{KR}{Recipe: Predict OpenMP Behaviour}
    \paragraph{Read directives} replicate vs worksharing.
    \paragraph{Scope each var} shared / private / reduction.
    \paragraph{Find races} a shared var written by $>1$ iteration $\to$ indeterminate.
    \paragraph{Count} multiply by thread count only in replicated regions.
\end{KR}

\begin{example2}{Ex.\ LZ2a (s7) How many prints?}
\begin{lstlisting}[language=C, style=basesmol, aboveskip=-0.05\baselineskip, belowskip=-0.5\baselineskip]
#pragma omp parallel              // (A)
for (int i=0;i<n;i++) printf("Hi\n");
#pragma omp parallel
#pragma omp for                   // (B)
for (int i=0;i<n;i++) printf("Hi\n");
\end{lstlisting}
    With $T$ threads and $n=5$, how many lines each?
     \\ \solution{Solution: }
    (A) \important{$T\times5$} (loop replicated on every thread); (B) \important{$5$} (iterations shared once). \texttt{parallel} makes threads; \texttt{for} splits work.
\end{example2}

\begin{example2}{Ex.\ LZ2b (s8) Value of \texttt{b}?}
\begin{lstlisting}[language=C, style=basesmol, aboveskip=-0.05\baselineskip, belowskip=-0.5\baselineskip]
int a=7,b=0,n=10;
#pragma omp parallel for shared(a,b)
for(int i=0;i<n;i++) b=a+i;        // section 1
#pragma omp parallel for shared(a) private(b)
for(int i=0;i<n;i++) b=a+i;        // section 2
\end{lstlisting}
     \solution{Solution: }
    S1 \texttt{shared}: \important{race} $\to$ any of $7..16$. S2 \texttt{private}: copies discarded $\to$ global \texttt{b} unchanged. Final \texttt{b} = whatever S1 left.
\end{example2}

\begin{example2}{Ex.\ LZ2c (s9) Reduction}
    Sum \texttt{a[0..8]} with no race; how does \texttt{reduction} work?
     \\ \solution{Solution: }
    \texttt{\#pragma omp parallel for reduction(+:sum)}: each thread builds a \important{private partial} (e.g.\ \texttt{a[0..2]}, \texttt{a[3..5]}, \texttt{a[6..8]}); after the join OpenMP adds the partials.
\end{example2}

\subsection{LZ3 IPC Synchronisation (Critical Sections)}

\mult{2}

\begin{definition}{Critical Section}
    Code only one thread may enter at a time (serialised). Entry must be uninterruptable (\texttt{lock()}); careless use $\to$ deadlock/livelock.
\end{definition}

\begin{concept}{Mechanisms}
    \begin{itemize}
        \item ISA \emph{atomic} (test-and-set)
        \item \texttt{omp atomic} (one op) / \texttt{omp critical [name]} (block)
        \item \texttt{sem\_*} (processes), \texttt{pthread\_mutex\_*} (threads, owner unlocks)
        \item \texttt{omp\_*\_lock} (expensive)
    \end{itemize}
\end{concept}

\multend

\begin{KR}{Recipe: Pick a Sync Primitive}
    single memory op $\to$ \texttt{atomic}; a block $\to$ \texttt{critical}; accumulate $\to$ \texttt{reduction}; processes $\to$ semaphore; threads $\to$ mutex; reusable handle $\to$ omp lock.
\end{KR}

\begin{example2}{Ex.\ LZ3 Increment a shared counter}
    Give three ways to make \texttt{count++} safe across threads, and the cheapest.
     \\ \solution{Solution: }
    \texttt{\#pragma omp atomic}, \texttt{\#pragma omp critical}, or a mutex/semaphore. \important{Cheapest = \texttt{atomic}} (uses a HW test-and-set; \texttt{critical}/locks are heavier).
\end{example2}

\subsection{LZ4 Thread Synchronisation (Barriers)}

\mult{2}

\begin{concept}{Barriers}
    \begin{itemize}
        \item implicit at end of a parallel region
        \item \texttt{\#pragma omp barrier} = explicit wall
        \item \texttt{nowait} drops an implicit barrier
        \item \texttt{ordered} = run the marked part in loop order
    \end{itemize}
\end{concept}

\begin{definition}{Barrier Stages}
    \begin{itemize}
        \item arrival (trapping): wait-queue \emph{or} busy-loop
        \item release: all trapped threads let go at once
    \end{itemize}
\end{definition}

\multend

\begin{KR}{Recipe: Control Barriers}
    need a wall $\to$ barrier; want overlap (no wait) $\to$ \texttt{nowait}; need ordered output $\to$ \texttt{ordered}; tune spin-vs-block via \texttt{OMP\_WAIT\_POLICY}.
\end{KR}

\begin{example2}{Ex.\ LZ4 (s18) Barrier internals}
    On arrival, wait-queue or busy-loop which is preferable? Once released, who goes first?
     \\ \solution{Solution: }
    \important{Depends}: short wait + spare cores $\to$ busy-loop (spin); long/oversubscribed $\to$ wait-queue (block). Release order = \important{OS scheduler, nondeterministic}; use \texttt{ordered} if order matters.
\end{example2}

\subsection{LZ5 Scheduling \& Task Management}

\mult{2}

\begin{concept}{Task Management}
    \begin{itemize}
        \item \texttt{master}: only the master thread runs the block
        \item \texttt{single}: one arbitrary thread runs it (implicit barrier)
        \item \texttt{sections}/\texttt{section}: task parallelism each section is a task
    \end{itemize}
\end{concept}

\begin{definition}{Loop Schedules}
    \begin{itemize}
        \item \texttt{static[,chunk]}: assigned up-front (even loads)
        \item \texttt{dynamic[,chunk]}: grab from queue (uneven)
        \item \texttt{guided[,chunk]}: shrinking chunks
    \end{itemize}
    set via clause / \texttt{OMP\_SCHEDULE} / \texttt{omp\_set\_schedule}.
\end{definition}

\multend

\begin{KR}{Recipe: Choose a Schedule}
    even \& predictable $\to$ \texttt{static} (low overhead); uneven/unknown $\to$ \texttt{dynamic}; tapering work $\to$ \texttt{guided}. Use \texttt{sections} for distinct concurrent tasks.
\end{KR}

\begin{example2}{Ex.\ LZ5a (s22) Effect of \texttt{nowait}}
    What does \texttt{nowait} do on \texttt{\#pragma omp single}?
     \\ \solution{Solution: }
    \texttt{single} runs the block on one thread with an \emph{implicit barrier}. \texttt{nowait} \important{removes that barrier} $\to$ the other threads proceed to ``Do some calculations'' without waiting (safe only if no dependency).
\end{example2}

\begin{example2}{Ex.\ LZ5b Which schedule?}
    A 64-iteration loop on 4 threads where iteration cost varies widely. Which schedule?
     \\ \solution{Solution: }
    \important{\texttt{dynamic}} (or \texttt{guided}) threads grab work as they finish, balancing the uneven load. \texttt{static} would leave fast threads idle.
\end{example2}

\subsection{LZ6 Runtime Interface}

\begin{definition}{Thread-Count Precedence}
    \texttt{if} clause $\to$ \texttt{num\_threads} $\to$ \texttt{omp\_set\_num\_threads()} $\to$ \texttt{OMP\_NUM\_THREADS} $\to$ default ($\approx$ \#CPUs). May exceed cores up to \texttt{OMP\_THREAD\_LIMIT}. Static $\to$ exact; dynamic $\to$ upper bound.
\end{definition}

\begin{KR}{Recipe: Determine the Thread Count}
    Walk the precedence list top-down; the first one set wins. Query at run with \texttt{omp\_get\_num\_threads} / \texttt{omp\_get\_thread\_num}.
\end{KR}

\begin{example2}{Ex.\ LZ6 Who wins?}
    \texttt{OMP\_NUM\_THREADS=8}, but the region is \texttt{\#pragma omp parallel num\_threads(4)}. How many threads?
     \\ \solution{Solution: }
    \important{4} \texttt{num\_threads} outranks the environment variable in the precedence list.
\end{example2}

\subsection{Exam FS23 Q4: Dot Product (C \& OpenMP)}

\begin{concept}{Scenario}
    $220\times220$ matrix; compute the dot product of every pair of row-vectors $\to$ output matrix. Prototype on a reduced $4\times4$.
\end{concept}

\begin{KR}{Recipe: Sequential Baseline then Parallelise}
    \paragraph{Baseline} nested loops + an inner MAC; comment; handle edges (bounds, overflow, symmetry).
    \paragraph{Parallelise} add \texttt{\#pragma omp parallel for} (\texttt{collapse} nested loops), keep the accumulator \emph{private}, pick a schedule.
\end{KR}

\begin{example2}{Exam FS23 Q4.1 MAC count (1P)}
    How many MAC operations in the reduced ($4\times4$) form?
     \\ \solution{Solution: }
    \important{$4\times4 = 16$} (per the answer key the $4\times4$ output of dot products).
\end{example2}

\begin{example2}{Exam FS23 Q4.2 Single-core C (6P)}
    Outline commented C for a single-core CPU with standard ALU.
     \\ \solution{Solution: }
\begin{lstlisting}[language=C, style=basesmol, aboveskip=-0.05\baselineskip, belowskip=-0.5\baselineskip]
// dot product of every row-vector pair, NxN -> NxN (N=4)
for (int i = 0; i < N; i++) {        // first vector (row i)
  for (int j = 0; j < N; j++) {      // second vector (row j)
    long sum = 0;                    // wide type avoids overflow
    for (int k = 0; k < N; k++)      // MAC over elements
      sum += in[i][k] * in[j][k];
    out[i][j] = sum;                 // symmetric: out[i][j]==out[j][i]
  }
}
\end{lstlisting}
\end{example2}

\begin{example2}{Exam FS23 Q4.3 OpenMP version (3P)}
    Enhance the code for a 4-core machine using OpenMP directives.
     \\ \solution{Solution: }
\begin{lstlisting}[language=C, style=basesmol, aboveskip=-0.05\baselineskip, belowskip=-0.5\baselineskip]
// split the pair-loops across 4 cores
#pragma omp parallel for collapse(2) schedule(static)
for (int i = 0; i < N; i++)
  for (int j = 0; j < N; j++) {
    long sum = 0;                    // private (declared inside)
    for (int k = 0; k < N; k++)
      sum += in[i][k] * in[j][k];
    out[i][j] = sum;
  }
\end{lstlisting}
\end{example2}
